%-*- coding: UTF-8 -*-
% gougu.tex
% 勾股定理
\documentclass[UTF8,12pt]{ctexart} %中文
\usepackage{graphicx} %图片处理宏
\usepackage{float} %浮动处理宏
% \usepackage{setspace} %行间距

\linespread{1.6}
\large


\title{杂谈勾股定理}
\author{张三}
\date{\today}

% 声明参考文献的格式
\bibliographystyle{plain}

% 以上称为：导言区（preamble）

\begin{document}
% 正文部分，空行无实意
\maketitle
% 生成标题

\begin{abstract}
  这是一篇关于勾股定理的小短文。
\end{abstract}

\tableofcontents
\newpage
\section{勾股定理在古代}

西方称勾股定理为毕达哥拉斯定理，将勾股定理的发现归功于公元前6世纪的毕达哥拉斯学派 \cite{Kline}。该学派得到了一个法则，可以求出可排成直角三角形三边的三元数组。毕达哥拉斯学派没有书面著作，该定理的严格表述和证明则见于欧几里德《几何原本》的命题47：“直角三角形斜边上的正方形等于两直角边上的两个正方形之和。”证明是用面积做的。

我国《周髀算经》载商高（约公元前12世纪）答周公问：
\begin{quote}
\zihao{-4}\kaishu 勾广三，股修四，径隅五。
\end{quote}
又载陈子（约公元前7--6世纪）答荣方问：
\begin{quote}
\zihao{-4}\kaishu 若求邪至日者，以日下为勾，日高为股，勾股各自乘，并而开方除之，得邪至日。
\end{quote}
都较古希腊更早。后者已经明确道出勾股定理的一般形式。图1是我国古代对勾股定理的一种证明 \cite{quanjing}。

\begin{figure}[ht]
  \centering
  \includegraphics[scale=0.6]{xuantu.pdf}
  \caption{宋赵爽在《周髀算经》注中作的弦图（仿制），该图给出了勾股定理的一个极具对称美的证明。}
  \label{fig:xuantu}
\end{figure}

\section{勾股定理的近代形式}

勾股定理可以用现代语音表述如下：

\newtheorem{thm}{定理}
\begin{thm}[勾股定理]
直角三角形斜边的平方等于两腰的平方和。
可以用符号语言表述为：设直接三角形 ABC，其中 $\angle{C} =90^\circ$，则有
\end{thm}
\begin{equation}
  AB^2=BC^2+AC^2
\end{equation}

满足式（1）的整数称为\emph{勾股数}。第1节所说毕达哥拉斯学派得到的三元数组就是勾股数。下表列出了一些较小的勾股数：
\\
\\
% H 不浮动（由 float 宏包提供）
\begin{table}[H]
  \begin{tabular}{|rrr|}
    \hline
    直角边 $a$ & 直角边 $b$ & 斜边 $c$ \\
    \hline
    3 & 4 & 5\\
    5 & 12 & 13\\
    \hline
  \end{tabular}
  \qquad
  ($a^2+b^2=c^2$)
\end{table}

% \section*{参考文献} %section*{} 无编号
% math 参考文献(*.bib)文件名
% BiBTaTex 专用于处理文档文献列表
% \addcontentsline{toc}{⟨level⟩}{⟨title⟩}
\nocite{Kline}
\bibliography{math}\addcontentsline{toc}{section}{参考文献}\tolerance=500
 
\end{document}
